\documentclass[11pt]{article}
\usepackage{explorer}
\newcommand{\ExplorerLanguage}{EN}
\newcommand{\ExplorerNumber}{02}
\newcommand{\ExplorerKind}{Lesson and method}
\begin{document}
\ExplorerTitle{High school · olympiad thinking}{The pigeonhole principle}
\ExplorerSection{Before you begin}
Euclidean division and elementary counting.\par
\ExplorerSection{Ideas and method}
If \(n + 1\) objects are placed in n boxes, at least one box contains two objects. Otherwise, each box would contain at most one, giving at most n objects.\par
More generally, if N objects occupy k boxes, one box contains at least ⌈\(\frac{N}{k}\)⌉. The ceiling ⌈x⌉ is the smallest integer greater than or equal to x.\par
Choosing the boxes is the key: remainders, birth months, small grid squares or classes of integers. Explain why every object belongs to one box.\par
\ExplorerSection{Worked example}
Among 13 people, two were born in the same month: the 12 months are the boxes. This does not guarantee two birthdays on the same day.\par
\ExplorerSection{Explore}
Choose six distinct integers from 1 to 10. Group them into pairs summing to 11 and find a complete pair.\par
\end{document}
